\chapter*{Recommendations}
\addcontentsline{toc}{chapter}{Recommendations}
\label{ch:recommendations}

{}\textbf{RECOMMENDATION}

Based on the study findings, several actions are recommended to support and provide for a sustainable urban development and ecological conservation in Makurdi, Otukpo, and similar tropical cities.

\textbf{Strengthening Urban Planning and Land-Use Regulation:}

Enhancing urban sustainability in Benue State requires stronger planning frameworks and more effective land-use regulation. To achieve this fit, a more stricter zoning policies should be implemented to limit the uncontrolled expansion into ecologically sensitive areas. By establishing urban growth boundaries, it will help prevent further loss of forest and agricultural land, while the integration of the ecological and geospatial data into development plans would ensure and provide a more evidence-based decision-making.

\textbf{Conservation and Restoration of Green Spaces:}

The act and manner of protecting and restoring natural vegetation is critical for maintaining ecological integrity within rapidly expanding cities. Remaining forest fragments, riparian buffers, and other natural patches should be legally protected to preserve biodiversity hotpots. Reforestation and afforestation programs in degraded peri-urban and urban zones can help recover lost habitats. In addition, establishing urban greenbelts and ecological corridors would improve habitat connectivity and support species movement across the landscape.

\textbf{Sustainable Urban Landscaping and Biodiversity Enhancement:}

It is adviseable to step up the urban landscaping practices to prioritize ecological sustainability and biodiversity conservation. The act and use of native plant species in landscaping should be encouraged to maintain local biodiversity and ecosystem processes. This fit can be achieved through the following; community-led initiatives, such as tree-planting campaigns, school greening programs, and urban gardens, can play an important role in improving vegetation cover. Furthermore, regulating indiscriminate vegetation clearing during construction is crucial to reducing avoidable biodiversity loss.

\textbf{Soil Quality Management:}

The maintenance of a healthy urban soils is essential for supporting vegetation growth and urban ecosystem services. It can be achieved through proper soil monitoring programs which should be established to track changes in pH, organic matter, and nutrient status across different land-use zones. On the other hand, the practices such as composting, mulching, and the application of organic amendments can improve soil structure and fertility. Policies must also be enforced to minimize soil contamination arising from improper waste disposal and construction activities.

\textbf{Climate and Ecosystem Resilience Strategies:}

Urban resilience should be encouraged and enhanced by integrating the climate responsive and ecosystem-based strategies into city planning. Expanding urban tree cover helps regulate microclimate, mitigate heat-island effects, and enhance air quality. The Nature based solutions such as green roofs, permeable pavements, and wetland restoration should be incorporated into urban infrastructure to manage stormwater and improve ecological functioning. Additionally, promoting alternatives to fuelwood use can help reduce pressure on remaining vegetation.

\textbf{Community Engagement and Environmental Education:}

To gainfully achieve a long-term environmental sustainability, it requires the active community involvement and also the increased public awareness. Educational campaigns should be developed to inform residents about the ecological impacts of urbanization and the importance of biodiversity conservation. Government should always empower the local communities to participate in conservation planning and environmental stewardship strengthens collective responsibility for ecological resources. Incorporating environmental education into school curricula will also help nurture future generations of sustainability-minded citizens.

\textbf{Policy Implications}

This study findings carry significant policy implications for urban planning, biodiversity conservation, and environmental management in Benue State and similar rapidly growing tropical cities. The revealed 517\% increase in built-up area expansion as observed between a 34 year period, the extensive loss of forest and cropland, declining species diversity, soil chemical alteration, and reduced vegetation productivity collectively demonstrate that current patterns of urbanization are ecologically unsustainable. These findings have highlight the urgent need for integrated land-use policies that recognize vegetation and biodiversity as critical urban infrastructure rather than expendable land resources. The high decline in the plant diversity across the urbanization gradients have underscores the need for state-level policies that can prioritize the protection of natural habitats and enforce compliance with environmental regulations during urban expansion projects.

The strong negative influence of population density on vegetation productivity as observed and revealed from the the suggests that demographic pressures must be a central consideration in urban policy formulation. Planners should incorporate ecological carrying capacity assessments into development approvals to prevent overcrowding and reduce vegetation removal associated with informal settlement growth. It is also important to note the significance changes in soil pH and nutrient conditions also point to the need for policies promoting sustainable soil management, including restrictions on indiscriminate waste disposal, erosion control measures, and the incorporation of green infrastructure in housing and commercial developments. These measures would help maintain soil health, which is essential for sustaining urban vegetation, mitigating heat island effects, and enhancing ecosystem resilience.

Also, the observed reduction in NDVI values and EVI values across urbanization gradients (zones) have implied that the vegetation cover is rapidly diminishing, thereby compromising essential ecosystem services such as carbon sequestration, microclimate regulation, stormwater retention, and air purification. Policymakers should therefore prioritize the expansion of urban green spaces, including parks, community forests, riparian buffers, and tree-lined streets. Another way to out on this issues, is the integration of the nature-based solutions such as green roofs, permeable pavements, and restored wetlands into urban infrastructure could help address the environmental consequences of rapid urbanization while improving the aesthetic and ecological value of urban neighborhoods.

On the other hand, the broader strategic level, the results highlight the need for evidence-based urban planning that incorporates geospatial data, ecological indicators, and predictive modelling. There should be a synergy between the Urban planners and decision makers in Benue State and they could adopt geospatial tools to monitor land-use trends, identify ecological hotspots requiring protection, and forecast the ecological impacts of proposed development. In addition, the study findings provide a more and reliable scientific foundation for updating state and local environmental policies, including Environmental Impact Assessment (EIA) procedures, urban forestry strategies, and climate adaptation frameworks. One important aspect to ensure is the strengthening of the institutional collaboration among environmental agencies, town planning authorities, and the community stakeholders is essential to ensure that ecological considerations remain central to development decisions.

Overall, the ecological trends revealed in this study call for a holistic policy shift from reactive environmental management toward proactive, sustainability-cantered urban governance. Lastly, the integration of the biodiversity conservation sector, the soil protection sector, and green infrastructure into urban planning, Benue State can mitigate the ecological costs of urbanization while promoting healthier, more resilient, and more liveable urban environments.

{}\textbf{Contribution to Knowledge}

The study highlighted numerous and important contributions to the scientific understanding of urban ecology in tropical environments. First, it introduces an integrated assessment framework that combines remote sensing, field-based biodiversity surveys, soil analyses, and statistical modelling, offering a methodological blueprint for similar research in other tropical cities. The second part, is that it provides a rare empirical evidence from a poorly studied region by delivering one of the few comprehensive ecological assessments of urbanization impacts in Benue State, thereby expanding the limited literature on West African urban ecosystems. The third part is that, the study demonstrates clear and significant urban gradient effects, showing how plant diversity, soil properties, and vegetation productivity change across rural, peri-urban, and urban zones, and thus contributing to broader global insights on ecological transitions in developing-world cities. Then, the fourth part is that, the identification of population density as the strongest predictor of vegetation decline advances understanding of the demographic pressures that drive ecosystem degradation in tropical urban landscapes. And finally is that, the study offers policy-relevant ecological insights by presenting biologically grounded recommendations for sustainable urban planning, biodiversity conservation, and ecosystem restoration, directly enriching applied urban ecology and environmental management.

{}\textbf{List of publications on the topic of the thesis}

Articles in journals included in the List of peer-reviewed scientific publications, in which the main scientific results of dissertations for the degree of Candidate of Sciences, for the degree of Doctor of Sciences should be published:

1. Shaibu, O., Omokaro, G.O., Onyamoche, E.P., Tafese, G.T., Nafula, Z.S., Niambe, O.K., Nnoli, E.C. (2025). Assessment of urbanization\textquotesingle s environmental impact in Makurdi Metropolis, Benue State, Nigeria. International Journal of Environmental Impacts, Vol. 8, No. 2, pp. 393-400. \url{https://doi.org/10.18280/ijei.080218} ( Scopus Q3 and WOS Q3)

2. Shaibu ochoche, Pinaev Vladimir Evgenievich., Emaikwu Patience Onyamoche,Kavhiza Nyasha.John, Opara Onyinyechi Grace, Niambe Obed Kohol. (2025). A social approach to assessing the impact of urbanization on natural resources: A case study of Makurdi-Benue State, Nigeria International Research Journal №6 (156). --- URL: \url{https://research-journal.org/en/archive/6-156-2025-june/10.60797/IRJ.2025.156.82} \url{https://doi.org/10.60797/IRJ.2025.156.82} Issue: № 6 (156), 2025 (VAK K1(SCOPUS Q2, WOS Q1)

3. Emmanuel Igwe,Jeremiah Akomaye Ikwen' Shaibu Ochoche, Aloye Racheal Aniah' Joseph Ikwun Agiopu Ogar.(2025) Roles of Artificial intelligence in predictive Agroecological modelling and sustainable land management in response to climate viability in Moscow Region. XXVth International Multidisciplinary Scientific GeoConference Surveying, Geology and Mining, Ecology and Management (SGEM) DOI~\textbf{10.5593/sgem2025/2.1/s07.07}~ISSN~\textbf{1314-2704}~ISBN~\textbf{978-619-7603-89-7} Conference Proceeding awaiting publishing Q2 SCOPUS

4. Shaibu Ochoche. (2025) Urbanization Impacts on Plant Diversity in Otukpo, Benue State, Nigeria. Young Scientific and Practical Conference. Department of Foreign Languages of the Faculty of Philology, Institute of Ecology RUDN University. (RSCI)

5. Ochoche Shaibu, Obed Kohol Niambe, Godspower Oke Omokaro, Zipporah Simiyu Nafula. (2025). Socio-ecological consequences of urban growth in dutsen kura gwari, Niger state, Nigeria, ``Международный научно-исследовательский центр ``Endless Light in Science'' Conference proceeding, DOI 10.24412/2709-1201-2025-31-93-101 (VAK, Library.ru, google scholar)

6. Godspower Oke Omokaro, Ikioukenigha Michael, Paul Obokparo Ewubare , Zipporah Simiyu Nafula , Obed Kohol Niambe , Shaibu Ochoche, Abdulmajeed Allan Chikukula and Elohor-Oghene Amarie. (2025). Soil fertility management practices among smallholder farmers in Ughelli North, Delta State, Nigeria. (2025). Tropical Agriculture, 102(3), 358--367. \url{https://journals.sta.uwi.edu/ojs/index.php/ta/article/view/9040}

(SCOPUS Q4 \& WOS Q4)

7. Shaibu Ochoche, Emmanuel Igwe, Vladimir Evgenievich Pinaev Emaikwu Patience Onyamoche (2025). Spatiotemporal Analysis of Land Use and Land Cover Changes Driven by Urbanization in Makurdi, Nigeria (1990 to 2024): Implications for Cropland Loss, Environmental Risk, and Sustainable Urban Development (SCOPUS Q3 WOS Q3) UNDER REVIEW

